\setchapterpreamble[u]{\margintoc}
\chapter{Derivatives and the Chain Rule}
\labch{derivatives}

\sources{Hass, Heil and Weir, \emph{Thomas' Calculus}, 14th ed.\ \cite{thomas2018calculus},
Chapters~2--3; MIT 18.01 \cite{ocw1801}.}

Part~II covers only the calculus that optimisation needs. That means
differentiation, and in particular the chain rule, which is the entire
mathematical content of training a network by gradients. Integration appears
only where \refch{expectation} needs it.

\section{The derivative as a local linear approximation}
\labsec{derivative-definition}

\begin{definition}
For \(f:\R\to\R\), the derivative at \(a\) is
\[
	f'(a)=\lim_{h\to0}\frac{f(a+h)-f(a)}{h},
\]
when the limit exists.
\end{definition}

The definition is a statement about approximation. Rearranged, it says
\[
	f(a+h)=f(a)+f'(a)h+o(h),
\]
where the remainder shrinks faster than \(h\). In words: near \(a\), the function
is the straight line \(f(a)+f'(a)h\), plus an error too small to matter at first
order.

\marginnote{This reading is the one that generalises. In \refch{gradients} the
straight line becomes a plane and \(f'(a)\) becomes a vector; in
\refch{jacobian-hessian} the plane becomes an affine map and the vector becomes
a matrix. The idea never changes.}

Differentiability implies continuity, but not conversely: \(f(x)=|x|\) is
continuous at \(0\) and has no derivative there, because the one-sided limits
disagree. This matters in practice, since the rectified linear unit
\(\max(0,x)\) has exactly this kink and is handled by convention rather than by
theorem.

\section{The rules}
\labsec{diff-rules}

For differentiable \(f\) and \(g\),
\[
	(f+g)'=f'+g',
	\qquad
	(fg)'=f'g+fg',
	\qquad
	\Bigl(\frac{f}{g}\Bigr)'=\frac{f'g-fg'}{g^{2}},
\]
together with \((x^{n})'=nx^{n-1}\), \((e^{x})'=e^{x}\), and
\((\ln x)'=1/x\) for \(x>0\).

\begin{example}
The logistic function \(\sigma(z)=\dfrac{1}{1+e^{-z}}\) satisfies
\[
	\sigma'(z)=\frac{e^{-z}}{(1+e^{-z})^{2}}=\sigma(z)\bigl(1-\sigma(z)\bigr).
\]
The derivative is expressible in terms of the value itself, which is why
implementations cache \(\sigma(z)\) during the forward pass and reuse it going
backwards.
\end{example}

\section{The chain rule}
\labsec{chain-rule}

\begin{theorem}
If \(g\) is differentiable at \(a\) and \(f\) is differentiable at \(g(a)\), then
\(f\circ g\) is differentiable at \(a\) and
\[
	(f\circ g)'(a)=f'\bigl(g(a)\bigr)\,g'(a).
\]
\end{theorem}

The linear-approximation reading makes the statement almost obvious: \(g\)
multiplies a small input change by \(g'(a)\), then \(f\) multiplies the result by
\(f'(g(a))\), so the composed effect is the product of the two factors.

\begin{example}
For \(f(x)=\sigma(wx+b)\) with \(w,b\) fixed,
\(f'(x)=\sigma'(wx+b)\cdot w\). Differentiating instead with respect to the
\emph{parameter} \(w\) gives \(\sigma'(wx+b)\cdot x\). Same rule, different
variable held fixed --- and it is the second version that training uses.
\end{example}

\marginnote{Almost every gradient in a neural network is this calculation
repeated. What makes networks distinctive is not the rule but the number of
times it is applied and the care taken over the order of the arithmetic.}

\section{Long compositions}
\labsec{long-compositions}

A network of \(L\) layers is a composition
\(f=f_L\circ f_{L-1}\circ\cdots\circ f_1\), and the chain rule gives
\[
	f'(x)=f_L'\bigl(z_{L-1}\bigr)\cdot f_{L-1}'\bigl(z_{L-2}\bigr)\cdots f_1'(x),
\]
where \(z_k\) is the output of layer \(k\). Two consequences deserve naming even
in the scalar case.

First, the derivative is a \emph{product} of \(L\) factors. If the factors are
systematically smaller than one the product decays geometrically; if
systematically larger it explodes. This is the vanishing and exploding gradient
problem, visible here with no matrices involved.

Second, the intermediate values \(z_k\) are needed to evaluate the factors. They
are computed on the way forward and must be kept, which is why memory during
training scales with depth.

\begin{remark}
Multiplication is associative but not equally cheap in every order. Accumulating
the product from the output end backwards keeps every intermediate quantity the
same shape as the output; accumulating from the input end does not.
\refsec{backprop} makes this precise once the objects are matrices.
\end{remark}

\section{What is not covered}
\labsec{calculus-scope}

Techniques of integration, infinite series, and differential equations are
outside the scope agreed for these notes. Definite integrals appear in
\refch{expectation} only to the extent that expectations of continuous random
variables require them, and the fundamental theorem of calculus is used there
without a separate development.
